This thesis examined whether the wording and timing of Telegram pump-and-dump signals help explain the immediate price reaction after the target coin is revealed. The answer is qualified, but clear. The public signal can matter, especially when urgency appears in a market setting where coordinated demand can still move the price. At the same time, the results do not support a general claim that more intense hype, more urgent wording, or longer instructions always produce higher returns. The effect depends on both the communication design of the pump and the market conditions around the targeted asset.

The research problem is important because pump-and-dump schemes create artificial price movements that are not based on a real change in the value of the asset. In these events, the message is not ordinary financial news. It is a tool for making many traders act at the same time. This makes the style of the signal relevant for understanding how manipulation works in fast and speculative cryptocurrency markets.

The study was built around its main hypotheses, stated here in full terms. The first hypothesis expected stronger values of the retained signal-characteristic vector, namely \textit{Hype Intensity}, \textit{Urgency Index}, and \textit{Instruction Density}, to have positive estimated effects on short-term cryptocurrency returns after the coin reveal. The second hypothesis expected the estimated effects of \textit{Hype Intensity}, \textit{Urgency Index}, and \textit{Instruction Density} to be stronger in lower-cap assets than in higher-cap assets, and to weaken as \textit{Market Capitalization} increases. The third hypothesis expected observable market patterns up to the reveal, combined with the same signal characteristics, to contain information that can help classify the binary \textit{Success}/\textit{Non-Success} outcome.

The empirical analysis used the complete pump-event sample constructed in the thesis. The data combined Telegram messages, exchange OHLCV market records, and historical CoinMarketCap metadata. The main outcome was the \textit{10-Minute Logarithmic Return} after the reveal. The causal part of the thesis used DML to estimate the relationship between signal characteristics and short-term returns after adjusting for pre-reveal market conditions, asset characteristics, and channel differences. The analysis also compared results across \textit{Market Capitalization} strata. A separate predictive part used XGBoost to test whether the binary \textit{Success}/\textit{Non-Success} outcome, defined by positive versus non-positive \textit{10-Minute Logarithmic Return}, could be classified from information available up to the reveal.

The first hypothesis receives partial support, mainly through \textit{Urgency Index}. Among the three signal characteristics, \textit{Urgency Index} has the clearest positive estimated role in the small-cap segment. \textit{Hype Intensity} and \textit{Instruction Density} show weaker and less stable patterns. The conclusion is therefore not that more intense hype, more urgent wording, or longer instructions always raise prices. A safer reading is that signal characteristics can matter, with urgency appearing most clearly when the target market is small enough to react, but still active enough for the reaction to be measured.

The second hypothesis receives the clearest but still qualified support. Market size changes the estimated effects of signal characteristics, especially \textit{Urgency Index}, which means that market structure is not only background information. It changes how the signal can translate into a price response. At the same time, the evidence does not support a simple rule that the smallest assets are always the easiest to move. The result is better understood as a conditional market-size pattern.

The third hypothesis is supported as a forecasting result. The \textit{Channel ID}-agnostic XGBoost model shows useful holdout performance when classifying binary positive versus non-positive short-window outcomes on unseen channels. This result is important, but it has to remain separate from the causal findings. Prediction answers a different question from causal explanation, and useful predictive features are not necessarily the direct causes of the return.

The thesis contributes to the study of cryptocurrency manipulation by treating pump messages as part of the manipulation mechanism, not only as labels for abnormal price movements. The public message is the moment where coordination becomes visible, and the empirical results suggest that this communication layer can be connected to market outcomes within clear limits. The contribution is also methodological, because the thesis keeps causal estimation and prediction as related but separate tasks.

These conclusions should be read within the limits of the research design. The study is observational, so unmeasured confounding cannot be fully ruled out. DML reduces this concern by adjusting for available pre-reveal information, but it cannot make the signals random. The \textit{10-Minute Logarithmic Return} captures the immediate market reaction, not full event profitability. It does not include the practical difficulty of entry and exit, transaction fees, or slippage. \textit{Market Capitalization} is useful as a practical proxy for market depth, but it is not the same as direct order-book liquidity. The micro-cap estimates remain difficult to interpret because trading is sparse in the smallest assets, and channel or organizer patterns may still affect both signal style and outcomes.

Future research could test the same mechanism with richer order-book depth data, execution and slippage information, longer post-event windows, more recent market periods, and broader market coverage.

Overall, the thesis shows that understanding cryptocurrency pump-and-dump schemes requires attention to both sides of the process. The social design of the signal helps organize traders at the moment of the reveal, while market conditions shape whether coordinated demand can move the price.
